\chapter{Resultados} \label{cap:05}

% Caso sejam somente resultados parciais, o nome do capítulo deve ser "Resultados Preliminares".
% Para a versão final, o nome do capítulo deve ser "Resultados e Discussão".

Nesta seção são apresentados os resultados obtidos durante o desenvolvimento da pesquisa. Os resultados podem incluir diferentes tipos de artefatos, tais como textos produzidos, diagramas, documentação de sistemas, interfaces desenvolvidas, além de análises provenientes dos dados coletados.

Ressalta-se que tanto resultados positivos quanto negativos devem ser apresentados, uma vez que contribuem para a construção do conhecimento científico e auxiliam futuros pesquisadores na compreensão das abordagens adotadas.

Dentre os resultados obtidos, destacam-se a implementação de funcionalidades do sistema e a validação de conceitos por meio de protótipos e testes.

\section{Exemplo de Implementação em JavaScript}

O trecho de código a seguir apresenta uma função responsável por consumir uma API e exibir os dados no console.

\begin{lstlisting}[caption={Requisição a uma API utilizando JavaScript}]
async function buscarProdutos() {
  try {
    // Faz requisicao para a API
    const resposta = await fetch('https://api.exemplo.com/produtos');
    
    // Converte resposta para JSON
    const dados = await resposta.json();

    // Percorre os produtos retornados
    dados.forEach(produto => {
      console.log(`Produto: ${produto.nome} - Preco: R$ ${produto.preco}`);
    });

  } catch (erro) {
    console.error('Erro ao buscar produtos:', erro);
  }
}

// Executa a funcao
buscarProdutos();
\end{lstlisting}

O código apresentado demonstra a utilização de programação assíncrona com \textit{async/await}, permitindo a comunicação com serviços externos por meio de APIs. Esse tipo de implementação foi fundamental para a integração entre o sistema desenvolvido e as fontes de dados utilizadas no projeto.

Além disso, outros resultados incluem a criação de interfaces, testes realizados e validação das funcionalidades propostas, os quais podem ser apresentados por meio de figuras, tabelas e gráficos ao longo desta seção.
